% Generated draft: requires mathematical and source review.
\documentclass[11pt]{article}
\usepackage[margin=25mm]{geometry}
\usepackage{fontspec}
\setmainfont{Latin Modern Roman}
\usepackage{amsmath,amssymb,mathtools,unicode-math}
\setmathfont{Latin Modern Math}
\usepackage{enumitem,needspace,xurl,hyperref}
\setlength{\parindent}{0pt}
\setlength{\parskip}{5pt}
\setlength{\emergencystretch}{3em}
\newcommand{\N}{\mathbb{N}}
\newcommand{\Z}{\mathbb{Z}}
\newcommand{\Q}{\mathbb{Q}}
\newcommand{\R}{\mathbb{R}}
\newcommand{\C}{\mathbb{C}}
\newcommand{\F}{\mathbb{F}}
\newcommand{\A}{\mathbb{A}}
\newcommand{\PP}{\mathbb{P}}
\newcommand{\E}{\mathbb{E}}
\newcommand{\eps}{\varepsilon}
\newcommand{\dd}{\,\mathrm{d}}
\begin{document}
\section*{UnsolvedMath ID: 800274}
\textbf{Source problem code:} catalogue entry
\subsection{Definitions and mathematical statement}
The source asks about holomorphic self-maps of the complex projective plane and leaves the degree range, the meaning of ``attractive basins,'' and the measure convention unspecified. The following is a precise formulation under explicit standard conventions; these conventions require review against the intended source meaning.

For a fixed positive integer \(d\), interpret degree as algebraic degree and put
\[
H_d=\left\{[F_0:F_1:F_2]: F_j\in\C[z_0,z_1,z_2]\text{ are homogeneous of degree }d,\quad \bigcap_{j=0}^2\{F_j=0\}\cap\PP^2(\C)=\varnothing\right\}.
\]
Here triples are identified under multiplication by a common nonzero scalar, and the corresponding map is
\[
[z_0:z_1:z_2]\longmapsto[F_0(z):F_1(z):F_2(z)].
\]
Thus \(H_d\) is an open subset of the projective space of coefficient triples. Interpret measure zero using the smooth measure class on this parameter space: a subset is null if its intersection with every projective coefficient chart has Lebesgue measure zero in the underlying real coordinates.

A periodic orbit \(\mathcal O=\{p,f(p),\ldots,f^{m-1}(p)\}\), where \(p\) has least period \(m\geq1\), is attracting if every eigenvalue of
\[
D(f^m)_p:T_p\PP^2(\C)\longrightarrow T_p\PP^2(\C)
\]
has modulus strictly less than \(1\). Fix any smooth Riemannian distance on \(\PP^2(\C)\) and define its full attracting basin by
\[
B_f(\mathcal O)=\left\{z\in\PP^2(\C):\lim_{n\to\infty}\operatorname{dist}(f^n(z),\mathcal O)=0\right\}.
\]
Under the convention that basins are counted by their distinct attracting periodic orbits, set
\[
E_d=\{f\in H_d: f\text{ has infinitely many distinct attracting periodic orbits}\}.
\]
The question is whether \(E_d\) has measure zero in \(H_d\), for each fixed degree \(d\) in the intended range.

The source does not specify whether the intended range is \(d\geq1\), \(d\geq2\), or another range. It also does not specify whether ``basins'' means the full sets \(B_f(\mathcal O)\), their connected components, immediate basins, or a broader notion of attraction. Consequently, the displayed formulation is conditional on the stated conventions and does not resolve these source ambiguities.
\subsection{Short English statement}
For each fixed intended degree, is the set of holomorphic self-maps of \(\PP^2(\C)\) with infinitely many attracting basins null in coefficient parameter space? The formulation above counts basins by distinct attracting periodic orbits and uses the smooth measure class; the source does not explicitly fix these conventions or the degree range.
\subsection{Sources}
\begin{itemize}
\item Local UnsolvedMath registry, permanent ID \(800274\), exact source code \texttt{catalogue entry}, title ``Several Complex Variables Problem List --- Problem 73.'' The supplied registry attributes the problem to Buzzard, Fornaess, and Krantz, with proposed year 1995, and labels it \texttt{open}. This is a supplied status label, not an independently verified assessment.
\item Buzzard, Fornaess, Krantz, \emph{The Several Complex Variables Problem List} (1995), Problem 73. Supplied source URL: \url{https://arxiv.org/abs/math/9509202}. The background records access on 2026-10-16 and extraction by \texttt{arxiv\_latex}; these are supplied provenance statements, not independent source verification for this draft.
\end{itemize}
\end{document}
